% ============================================================
% CAPÍTULO 6 — DOCUMENTAÇÃO DA UTILIZAÇÃO DE IA GENERATIVA
% Sistema de Gestão Integrada — Clube Amigos de Polvoreira (CAP)
% ============================================================

\chapter{Documentação da Utilização de Inteligência Artificial Generativa}\label{chap:documentacao_ia}

A utilização de modelos de linguagem neste projeto foi uma parte central do nosso processo de trabalho, presente em praticamente todas as fases do desenvolvimento. Este capítulo tem dois objetivos: registar de forma transparente como e onde usámos IA generativa ao longo do projeto, e refletir sobre o que correu bem, o que correu mal e como a equipa validou o que vinha dos modelos.

O capítulo está organizado por fase do projeto. Quando houve interações relevantes, os \textit{prompts} utilizados aparecem nas caixas azuis com o modelo usado e o resultado obtido. Toda a documentação de IA está concentrada aqui num único capítulo.

% ============================================================
% SECÇÃO 6.1 — ABORDAGEM GERAL E FERRAMENTAS
% ============================================================

\section{Abordagem Geral e Ferramentas Utilizadas}\label{sec:ia_abordagem}

Não é possível falar de como usámos os LLMs sem primeiro identificar os modelos que utilizámos e a forma como interagimos com eles. Esta secção apresenta esse enquadramento.

\subsection{Modelos Utilizados}

Ao longo do projeto experimentaram-se vários modelos, de diferentes fornecedores, mas dois deles destacaram-se claramente na frequência de utilização:

\begin{itemize}
    \item \textbf{Google Gemini} (família Gemini 2.5 Pro e Flash), utilizado maioritariamente pela interface Web. Foi o nosso modelo de referência para tarefas exploratórias: geração de documentos do clube, estruturação do questionário, redação de entrevistas simuladas e \textit{brainstorming} de opções arquiteturais. A vantagem do Gemini Web é a rapidez de iteração --- responde em segundos, suporta \textit{upload} de ficheiros de contexto e mantém o histórico da conversa numa janela acessível.

    \item \textbf{Anthropic Claude} (modelos Sonnet 4.5 e 4.6), utilizado quase sempre através da interface de terminal (Claude Code). Foi o modelo preferido para tarefas que envolvessem múltiplos ficheiros do projeto em simultâneo: redação em bloco do detalhamento dos requisitos, geração de código C\#, refactoring transversal e correções gramaticais aplicadas ao relatório inteiro.
\end{itemize}

Em fases mais específicas, recorremos pontualmente a:

\begin{itemize}
    \item \textbf{Vercel v0}, utilizado na fase de \textit{design} de interfaces para gerar protótipos das principais páginas da aplicação (\textit{dashboards} dos quatro perfis de utilizador) a partir de descrições textuais. A diferença para uma resposta normal de LLM é que o \textit{v0} devolve diretamente código React/JSX funcional, que serviu de base para a implementação da camada de apresentação.

    \item Modelos da OpenAI (GPT-4o e GPT-5), pontualmente, sobretudo para comparar respostas em decisões críticas e para confirmar diagnósticos quando o Gemini ou o Claude pareciam estar a alucinar.
\end{itemize}

A combinação Gemini + Claude foi a que melhor funcionou para nós. O Gemini é bom para explorar ideias e fazer perguntas rápidas na Web; o Claude no terminal é melhor para trabalhar diretamente nos ficheiros do projeto.

\subsection{Estratégias de Interação: Web vs CLI}

Em termos práticos, usámos os modelos de duas formas distintas, com perfis de utilização muito diferentes.

\paragraph{Interface Web (oficial dos fornecedores).}
A versão Web (chat.google.com/gemini, claude.ai) é prática para resolver problemas rápidos, fazer perguntas pontuais, pedir correções e obter opiniões breves. Foi a estratégia dominante na fase de levantamento de requisitos: geração dos documentos do clube (atestado médico, relatório de jogo, fatura de quota), simulação de entrevistas com os \textit{stakeholders} e estruturação do questionário online. A vantagem é a velocidade de iteração e a possibilidade de carregar facilmente ficheiros de contexto. A limitação é que o modelo não tem acesso direto aos ficheiros do projeto: cada interação começa praticamente do zero, e cabe a quem usa colar o contexto necessário.

\paragraph{Interface de terminal (Claude Code, Gemini CLI).}
A segunda estratégia consiste em executar o LLM diretamente no diretório do projeto. É mais custosa em termos de tempo de preparação e \textit{tokens} consumidos, mas oferece vantagens decisivas:

\begin{itemize}
    \item \textbf{Acesso direto aos ficheiros do trabalho.} O modelo lê e edita diretamente os ficheiros \texttt{.tex}, \texttt{.cs}, \texttt{.puml} etc. Isso elimina o esforço manual de colar código, copiar respostas e reformatar — particularmente útil para tarefas repetitivas, como corrigir erros gramaticais ao longo de todo o relatório, uniformizar formatação de imagens ou alinhar IDs de requisitos entre capítulos.
    \item \textbf{Contexto persistente do projeto.} Como o modelo vê toda a estrutura de diretórios, as respostas são significativamente mais precisas. Há menos alucinações: o modelo sabe que classes existem, que ficheiros estão presentes e que decisões já foram tomadas.
    \item \textbf{Operações em \textit{batch}.} Tarefas que seriam tediosas de fazer caixa-a-caixa na Web (por exemplo, gerar a tabela detalhada para 36 requisitos) podem ser delegadas num único \textit{prompt} bem construído.
\end{itemize}

O risco associado a esta estratégia é também claro: como o modelo opera diretamente sobre o ficheiro, é mais fácil deixar passar uma alteração indesejada. A equipa minimizou este risco através de duas práticas: \textit{commits} frequentes (para conseguir reverter alterações com um \texttt{git diff}/\texttt{git restore}) e revisão humana dos \textit{diffs} antes de qualquer alteração não trivial.

\paragraph{Combinação das duas estratégias.}
Ao usar as duas estratégias em paralelo, conseguimos tirar o melhor de cada uma sem depender totalmente de nenhuma. A Web servia para explorar ideias e fazer perguntas rápidas; o CLI para aplicar decisões e propagar correções pelo projeto todo.

\subsection{Utilização Passiva: LLM como Assistente de Escrita}

Uma utilização menos visível, mas importante de registar, é o uso passivo dos modelos como assistentes de escrita do próprio relatório. Não se trata de gerar conteúdo novo: trata-se de pedir ao modelo para corrigir gralhas, suavizar parágrafos demasiado robóticos, sugerir variações de vocabulário ou validar a sintaxe \LaTeX{}. Esta utilização foi sobretudo feita com Claude Code, dada a sua capacidade de operar sobre todo o documento de uma só vez.

Importa sublinhar uma decisão da equipa: nas situações em que se delegou ao modelo a redação de várias secções consecutivas (como aconteceu no detalhamento inicial dos requisitos), o texto produzido tendia a ficar uniforme demais e a perder naturalidade. Foi necessário voltar atrás e reescrever passagens para introduzir variação no estilo. Esta reflexão sobre o risco de uniformização está desenvolvida na Secção~\ref{sec:ia_validacao}.

% ============================================================
% SECÇÃO 6.2 — APLICAÇÃO POR FASE DO PROJETO
% ============================================================

\section{Aplicação por Fase do Projeto}\label{sec:ia_fases}

\subsection{Fase 1 --- Definição e Contextualização do Sistema}

Nesta primeira fase, a equipa decidiu \textbf{não usar IA generativa}. A contextualização do CAP, a relação com a comunidade de Polvoreira e o entendimento dos recursos disponíveis dependiam do nosso contacto direto com o clube. Usar um LLM aqui teria dado um texto genérico que não serviria de base para as fases seguintes.

O Capítulo~\ref{chap:especificacao_modelacao} foi assim integralmente redigido pela equipa, sem recurso a modelos.

\subsection{Fase 2 --- Elicitação e Análise de Requisitos}\label{sec:ia_requisitos}

Esta foi a fase do projeto em que a utilização de IA teve mais peso. Os artefactos produzidos com auxílio de LLMs incluem os documentos sintéticos do clube, as entrevistas simuladas, o questionário online, a lista inicial de requisitos e o detalhamento exaustivo de cada um deles. Cada uma destas atividades é descrita nas subsecções seguintes, com o respetivo \textit{prompt} e a estratégia adotada.

\subsubsection{Definição da Estratégia de Levantamento}

Antes de avançar para a recolha propriamente dita, recorremos ao Gemini Web para enquadrar metodologicamente o levantamento. Apresentámos o contexto do CAP e pedimos uma lista de métodos adequados a um ecossistema misto (atletas menores, gerência, treinadores, secretaria). A resposta serviu de catálogo: a equipa filtrou-a para os três métodos que faziam sentido dentro do prazo do projeto --- análise documental, entrevistas semiestruturadas e questionário online.

\caixaprompt{Gemini (Web)}{No contexto do CAP quero que me digas quais são os melhores métodos para o levantamento de requisitos para o nosso novo sistema de software.}{Para um ecossistema tão rico e diversificado como o do Clube Amigos de Polvoreira, recomendo uma abordagem mista (triangulação de métodos), inspirada nas práticas de Sommerville. Aqui estão os métodos mais indicados para o vosso caso:

\begin{itemize}
    \item \textbf{Análise Documental:} Fundamental para perceber ``o que existe''. Devem recolher os atestados médicos, relatórios de jogo em papel e os livros de recibos de quotas atuais para perceber exatamente que dados precisam de ser transpostos para o digital.

    \item \textbf{Entrevistas Direcionadas (Semiestruturadas):} Essencial para captar dores específicas. Recomendo entrevistar representantes chave: um elemento da direção (para a vertente financeira/SAF-T), um treinador (sobre convocatórias e lesões) e um atleta.

    \item \textbf{Questionários Online:} Sendo um clube com muitos associados, a melhor forma de validar problemas à escala (ex: dificuldades de comunicação em grupos de WhatsApp ou métodos de pagamento preferidos) é enviando um formulário via email para todos os atletas e encarregados de educação.

    \item \textbf{Observação Direta (Etnografia):} Passar uma tarde na secretaria do clube num dia de jogo ou de inscrições para observar os estrangulamentos do processo manual.

    \item \textbf{Focus Groups (Grupos de Foco):} Juntar numa sala treinadores de diferentes modalidades para debater como deve funcionar o sistema de reserva do pavilhão sem gerar conflitos.

    \item \textbf{Prototipagem Rápida:} Criar ecrãs falsos (\textit{mockups}) do sistema e mostrar à gerência para garantir que a usabilidade se adequa à literacia digital dos funcionários.
\end{itemize}

Recomendo que escolham 2 a 3 destes métodos para garantir uma boa cobertura sem esgotar o tempo útil do projeto.}

\subsubsection{Identificação dos \textit{Stakeholders}}

Com a estratégia definida, o passo seguinte foi identificar todos os \textit{stakeholders} relevantes. Carregámos no Gemini um documento com a contextualização do clube e pedimos uma listagem categorizada. O modelo identificou quase todos os perfis que já tínhamos em mente --- atletas, treinadores, gerência, secretaria, sócios não-praticantes, federações e Autoridade Tributária ---, mas \textbf{deixou de fora os Encarregados de Educação}, presumivelmente por não ter ainda informação suficiente sobre a estrutura etária dos atletas do CAP. A equipa acrescentou manualmente esse perfil, que veio a revelar-se um dos mais importantes do sistema (RF-02).

\caixaprompt{Gemini (Web)}{No contexto do CAP, quero que identifiques todos os stakeholders do clube.}{Para o Clube Amigos de Polvoreira (CAP), a identificação correta dos \textit{stakeholders} é fundamental para perceber quem irá interagir com o novo sistema e quem será afetado por ele. Tendo em conta a estrutura típica de um clube desportivo português, dividi os intervenientes em duas categorias principais:

\textbf{\textit{Stakeholders} Primários (Utilizadores Diretos):}
\begin{itemize}
    \item \textbf{Atletas:} Os praticantes das diversas modalidades do clube.
    \item \textbf{Treinadores:} Os responsáveis técnicos das equipas, que farão convocatórias e registo de lesões.
    \item \textbf{Encarregados de Educação:} Os responsáveis legais pelos atletas menores, fundamentais para a gestão de pagamentos e autorizações.
    \item \textbf{Secretaria:} Os funcionários administrativos que gerem inscrições, validação de atestados e processos diários.
    \item \textbf{Gerência:} Os órgãos de direção e administração que necessitam de uma visão global financeira e desportiva.
\end{itemize}

\textbf{\textit{Stakeholders} Secundários (Afetados Indiretamente):}
\begin{itemize}
    \item \textbf{Sócios não-praticantes:} Membros associados sem atividade desportiva ativa.
    \item \textbf{Fornecedores e parceiros:} Entidades externas com relações comerciais com o clube.
    \item \textbf{Autoridades desportivas:} Federações e associações reguladoras das competições.
    \item \textbf{Autoridade Tributária (AT):} Entidade reguladora fiscal, essencial para efeitos de conformidade e exportação do SAF-T.
\end{itemize}}

\subsubsection{Geração dos Documentos do Clube}

A análise documental exigia que tivéssemos exemplos representativos dos documentos manipulados pelo CAP --- atestado médico desportivo, relatório oficial de jogo e fatura/recibo de quota. Como esses documentos não estão publicamente disponíveis num formato útil, optámos por sintetizá-los com auxílio do Gemini.

A geração não foi imediata. Houve três iterações:

\begin{itemize}
    \item \textbf{Iteração 1 --- texto livre.} Pedimos ao Gemini para gerar o atestado em texto comum. O resultado foi correto no conteúdo mas inutilizável visualmente: um bloco de texto sem cabeçalho, sem campos identificados e sem layout.
    \item \textbf{Iteração 2 --- LaTeX.} Pedimos a versão em LaTeX. O modelo gerou código \texttt{tabular} válido, mas o resultado renderizado ficou rígido e pouco realista --- não parecia um documento oficial.
    \item \textbf{Iteração 3 --- HTML.} Por sugestão do próprio modelo, pedimos o documento em HTML com estilos CSS embebidos. Esta versão saiu praticamente como queríamos logo no primeiro pedido. A pré-visualização foi convertida em PDF e usada como ponto de partida para a extração de requisitos.
\end{itemize}

Os três documentos finais (atestado, relatório de jogo e fatura) encontram-se no diretório \texttt{Relatorio/docs/docsclube/}.

\caixaprompt{Gemini (Web)}{Quero que geres um atestado médico desportivo em HTML, tipo os que se usam para a prática federada em Portugal. Mete os campos normais (médico, vinheta, validade, aptidão, dados do atleta). Depois diz-me que requisitos é que consegues tirar desse documento para o nosso sistema.}{}

\caixaprompt{Gemini (Web)}{Agora faz o mesmo mas para um relatório de jogo. Mete os campos que normalmente se preenchem depois de um jogo (adversário, resultado, quem jogou, golos, cartões, substituições, observações do treinador). E extrai os requisitos também.}{}

\caixaprompt{Gemini (Web)}{Agora faz uma fatura de quota do clube, também em HTML. Tem de ter os campos todos para ficar conforme com o SAF-T (NIFs, série de faturação, IVA, etc.). E lista os requisitos.}{}

\subsubsection{Simulação de Entrevistas}

Para complementar a análise documental, simulámos entrevistas semiestruturadas a três perfis representativos: gerência, treinador e atleta (com perspetiva do encarregado de educação). Esta opção foi assumida claramente como uma simulação --- não foram entrevistas reais ---, e o Gemini foi instruído a manter um registo realista de problemas operacionais portugueses (papelada, WhatsApp, atestados caducados, conflitos de pavilhão).

O resultado foi satisfatório logo no primeiro \textit{prompt} e praticamente não exigiu intervenção. As três transcrições estão guardadas em \texttt{Relatorio/docs/interviews/}.

\caixaprompt{Gemini (Web)}{Quero que geres uma entrevista simulada entre alguém da nossa equipa e um membro da gerência do CAP. Foca nos problemas reais tipo quotas, conflitos de instalações, obrigações fiscais. Faz com que pareça natural e que o entrevistado conte situações concretas.}{}

\caixaprompt{Gemini (Web)}{Agora faz o mesmo mas com um treinador. Foca-te nos problemas do dia a dia: comunicação por WhatsApp, dificuldade em fazer convocatórias, não saber se os atletas têm atestado válido, esse tipo de coisas.}{}

\caixaprompt{Gemini (Web)}{E agora uma com um atleta (ou o encarregado de educação se for menor). Sobre presenças, pagamento de quotas, atestados, e se gostava de ter uma área pessoal com as estatísticas dele.}{}

\subsubsection{Estruturação do Questionário Online}

Para escalar a recolha, decidimos enviar um questionário a todos os membros do clube com endereço de email registado. Como o Gemini não consegue criar formulários no Google Forms diretamente, a estratégia adotada foi diferente: pedimos primeiro a \textbf{estrutura} do questionário (secções, perfis-alvo, tipos de pergunta) e depois, num segundo momento, pedimos as \textbf{perguntas uma a uma}, num formato pronto a colar no Google Forms.

Houve perguntas que a equipa rejeitou (por serem demasiado genéricas ou repetitivas) e outras que foram refinadas até ficarem claras. O questionário final ficou com aproximadamente vinte perguntas, distribuídas pelas quatro secções definidas pelo modelo.

\caixaprompt{Gemini (Web)}{Quero que geres a estrutura de um questionário Google Forms para ser enviado a todos os stakeholders, e depois vamos gerar as perguntas 1 prompt de cada vez.}{Excelente abordagem. Para garantir que o questionário capte as necessidades de todos os intervenientes do Clube Amigos de Polvoreira de forma organizada, recomendo dividir o Google Forms nas seguintes secções principais, utilizando a lógica de ramificação (\textit{branching}) do formulário:

\begin{itemize}
    \item \textbf{Secção 1: Perfil do Inquirido:} Pergunta inicial de seleção múltipla (Atleta, Encarregado de Educação, Treinador, Sócio ou Secretaria/Gerência) para direcionar o utilizador.
    \item \textbf{Secção 2: Avaliação Geral (Todos):} Questões sobre os canais de comunicação atuais (WhatsApp, email, presencial) e as maiores dores na obtenção de informação do clube.
    \item \textbf{Secção 3: Módulos Específicos (Ramificados):}
    \begin{itemize}
        \item \textit{Para Treinadores:} Dificuldades nas convocatórias, assiduidades e relatórios de jogo.
        \item \textit{Para Pais/Atletas:} Experiência com o pagamento de quotas, envio de atestados e acesso ao calendário de jogos.
    \end{itemize}
    \item \textbf{Secção 4: O Novo Sistema:} Pergunta de resposta aberta para sugestão de funcionalidades prioritárias na nova plataforma.
\end{itemize}}

\vspace{0.3cm}
\noindent\textit{\textbf{Nota da equipa:} Por motivos de concisão, optámos por não documentar aqui a totalidade das iterações seguintes. As perguntas exatas e as suas opções de resposta foram desenvolvidas iterativamente, com um \textit{prompt} dedicado a cada pergunta, e coladas no Google Forms diretamente após validação interna.}

\caixaprompt{Gemini (Web)}{Gera agora a primeira pergunta da secção de comunicação. Quero perceber se as pessoas estão satisfeitas com o WhatsApp e as chamadas ou se preferiam outra coisa. Diz-me que tipo de pergunta usar e dá-me as opções.}{}

\subsubsection{Extração dos Requisitos a partir dos Artefactos}

Reunidos os três documentos sintéticos, as três entrevistas e as 47 respostas (na simulação) do questionário, passámos à fase de extração de requisitos. Para evitar enviesamentos da conversa anterior, abrimos um \textbf{novo chat} no Gemini e enviámos cada artefacto separadamente, pedindo apenas a \textbf{lista de nomes} dos requisitos (em texto simples) para podermos consolidar.

O primeiro resultado tinha bastante ambiguidade --- alguns requisitos apareciam repetidos com nomes ligeiramente diferentes, outros eram demasiado vagos para serem implementáveis. Explicámos o problema ao modelo e pedimos uma versão limpa, sem ambiguidades. O resultado dessa segunda iteração foi satisfatório.

Por fim, pedimos ao Gemini para dividir a lista consolidada em \textbf{requisitos funcionais} e \textbf{requisitos não funcionais}, com uma justificação curta para cada classificação. Esta divisão deu origem aos 23 RF e 13 RNF que aparecem no Capítulo~\ref{chap:analise_requisitos}.

\caixaprompt{Gemini (Web)}{Estou a mandar-te os documentos do clube (atestado, relatório de jogo, fatura), as 3 entrevistas e as respostas do questionário. Quero que tires uma lista de todos os requisitos que encontras, só o nome curto de cada um, sem descrições. Depois vamos consolidar.}{}

\caixaprompt{Gemini (Web)}{Essa lista tem problemas, há requisitos repetidos com nomes diferentes e outros que são vagos demais. Limpa isso: junta os duplicados, reescreve os que são vagos para algo mais concreto e tira os que não servem para nada. Devolve a lista limpa em texto simples.}{}

\caixaprompt{Gemini (Web)}{Agora divide a lista em requisitos funcionais e não funcionais. Para cada um diz-me numa linha porque é que o classificas assim.}{}

\subsubsection{Detalhamento Exaustivo dos Requisitos (Claude Code)}

Esta foi uma das fases mais pesadas em termos de trabalho bruto. Cada um dos 36 requisitos (23 RF + 13 RNF) precisava de ter uma tabela detalhada com ID, nome, área, ator principal, descrição, prioridade, origem, analista, pré-condições, pós-condições e critérios de aceitação --- mais um parágrafo explicativo que justificasse a sua existência. Fazer isto manualmente seria penoso e propenso a inconsistências.

A estratégia foi a seguinte: começámos no Gemini Web, num novo chat, com a lista limpa de requisitos como contexto. Gerámos manualmente as primeiras quatro tabelas (uma por requisito), para fixar o formato visual e validar a estrutura. Depois passámos para Claude Code, no diretório do projeto, com instrução para preencher os restantes requisitos seguindo exatamente o mesmo formato e estilo. O Claude leu os ficheiros de entrevistas e o questionário diretamente para inferir a origem e o analista de cada requisito.

O processo de geração em \textit{batch} demorou aproximadamente uma hora de execução do modelo, seguidas de duas horas de revisão pela equipa, com pequenos ajustes em prioridades e critérios de aceitação. Sem este auxílio, estimamos que o trabalho equivalente teria consumido pelo menos uma semana de redação manual.

\caixaprompt{Claude Code (CLI)}{Olha para o ficheiro chapters/03\_analise\_requisitos.tex, já tem 4 tabelas preenchidas como exemplo. Quero que faças as tabelas dos restantes requisitos exatamente no mesmo formato. Vai buscar a origem de cada requisito às entrevistas em docs/interviews/. Antes de cada tabela mete um parágrafo a explicar o requisito. Não inventes dados que não estejam nos documentos.}{}

\subsection{Fase 3 --- Arquitetura e Modelação do Software}\label{sec:ia_arquitetura}

Concluído o levantamento, a fase de arquitetura e modelação envolveu decisões técnicas (qual o padrão arquitetural, qual o motor de base de dados, qual a linguagem de \textit{backend}) e a produção de uma grande quantidade de diagramas UML. O uso de IA foi feito de forma diferente em cada um destes momentos.

\subsubsection{Decisões Arquiteturais}

Antes da reunião decisória, recorremos ao Gemini Web para obter uma síntese estruturada das opções disponíveis em cada um dos quatro eixos: padrão arquitetural global, motor de base de dados, \textit{stack} de \textit{backend} e paradigma de API. Para cada eixo, o modelo apresentou três candidatos com prós e contras face às restrições específicas do CAP (RGPD, SAF-T, equipa reduzida, prazo apertado). A equipa pesou as opções na reunião decisória e tomou a decisão final em conjunto. Em todas as decisões a opção final foi a sugerida pelo modelo --- mas isso aconteceu porque a equipa concordou com o raciocínio apresentado, não por delegação automática.

\caixaprompt{Gemini (Web)}{Temos de decidir a arquitetura do CAP. Dá-me os 3 melhores candidatos para cada uma destas decisões, com prós e contras de cada um: (1) padrão arquitetural, (2) base de dados, (3) tecnologia de backend, (4) tipo de API. Tem em conta que somos uma equipa de 4 pessoas, temos 16 semanas, precisamos de RGPD, SAF-T e de cronjobs para verificar atestados.}{}

\caixaprompt{Gemini (Web)}{Revê a escolha do Monólito Modular para o CAP. Há algum risco que devamos antecipar, dado que o sistema lida com dados clínicos, financeiros e desportivos ao mesmo tempo? Sê específico.}{}

\subsubsection{Modelação Estrutural: Modelo de Domínio}

Para a modelação estrutural adotámos uma estratégia uniforme: gerar primeiro o esqueleto do diagrama em PlantUML (com Gemini CLI), refinar manualmente no editor PlantUML e finalmente exportar para PNG.

O modelo de domínio foi o diagrama mais difícil de obter. O Gemini insistia em produzir classes UML completas, com atributos e métodos, em vez de um modelo de domínio puro (apenas entidades e relações). Foram necessárias três iterações para convencer o modelo a remover os atributos e a focar-se apenas nas entidades conceituais e nas suas relações --- mantendo as classes de associação onde fizessem sentido (\textit{Atleta no Plantel}, \textit{Item na Convocatória}, etc.). A versão final foi ainda ajustada pela equipa para incluir a especialização \textit{Pagamento de Quota}/\textit{Fatura} a partir da abstração \textit{Transação}.

\caixaprompt{Gemini CLI}{Lê os requisitos em chapters/03\_analise\_requisitos.tex e gera o modelo de domínio em PlantUML. Atenção: é modelo de domínio, só entidades e relações, não metas atributos nem métodos. Usa classes de associação tipo Atleta x Equipa, Item x Convocatória, etc. Grava em images/modelo\_dominio.puml.}{}

\caixaprompt{Gemini CLI}{Ainda meteste atributos nas entidades (nome, data no Atleta e no Treino...). Tira tudo, é modelo de domínio, só entidades e relações. Não é diagrama de classes.}{}

\caixaprompt{Gemini CLI}{Faz com que Pagamento de Quota e Fatura sejam especializações de uma classe Transação, para o SAF-T. Atualiza o .puml.}{}

\subsubsection{Diagrama de Classes}

O diagrama de classes aprofunda o modelo de domínio com detalhe técnico: visibilidade dos atributos, tipos de dados, assinaturas de operações e distinções claras entre herança, composição e associação. Aqui o Gemini CLI conseguiu gerar uma primeira versão razoável a partir do modelo de domínio, mas foi necessário corrigir manualmente algumas escolhas: tornar \texttt{Utilizador} e \texttt{Funcionario} abstratas, mover o \texttt{nif} para a classe base e acrescentar o método \texttt{isApto()} em \texttt{Atleta} --- decisões que foram inclusive sugeridas pelo modelo numa interação posterior, depois de lhe pedirmos uma revisão.

\caixaprompt{Gemini CLI}{A partir do modelo\_dominio.puml, gera agora o diagrama de classes em PlantUML. Desta vez quero tudo: atributos com tipos, operações, visibilidade. Organiza por módulo (Utilizadores, Desportivo, Clínico, Financeiro, Infraestrutura). Grava em images/diagrama\_classes.puml.}{}

\caixaprompt{Gemini (Web)}{Revê o diagrama de classes do subsistema Utilizadores. A hierarquia Utilizador $\rightarrow$ Atleta/Treinador/EE/Funcionário/Secretária/Gerência está correta? Falta algum atributo ou método importante?}{}

\subsubsection{Diagramas de Componentes e Implantação}

Os diagramas de componentes e implantação foram gerados em PlantUML diretamente pelo Gemini CLI, com base na arquitetura de Monólito Modular já decidida. A primeira versão do diagrama de componentes apresentava o ``API Gateway'' a executar lógica de negócio, o que é um antipadrão (\textit{fat gateway}). Numa segunda iteração, o modelo foi instruído a limitar o gateway a roteamento, autenticação e autorização, deixando a lógica nos módulos.

\caixaprompt{Gemini CLI}{Gera o diagrama de componentes em PlantUML. A arquitetura é monólito modular com 7 módulos, API Gateway em ASP.NET Core, PostgreSQL, e integrações externas com pagamentos e email/SMS.}{}

\caixaprompt{Gemini CLI}{O API Gateway que geraste está a fazer coisas a mais, tem lógica de negócio lá dentro. Isso é um fat gateway. Corrige para que o gateway só faça routing, autenticação e RBAC, o resto fica nos módulos.}{}

\caixaprompt{Gemini CLI}{Agora gera o diagrama de implantação. Cliente é SPA + React Native, servidor é .NET 7 em Linux, base de dados PostgreSQL, serviços externos por HTTPS.}{}

\subsubsection{Máquinas de Estado}

Foram construídas seis máquinas de estado para as entidades com ciclo de vida crítico: Atleta, Quota, Convocatória, Atestado, Reserva de Instalação e Lesão. O Gemini CLI gerou as primeiras versões a partir do modelo de domínio e dos requisitos correspondentes. As máquinas mais complexas (Atleta e Convocatória) precisaram de iteração manual para representar corretamente as transições paralelas (atleta pode ficar inapto por motivo clínico, financeiro ou disciplinar, em qualquer momento).

\caixaprompt{Gemini CLI}{Gera as máquinas de estado em PlantUML para: Atleta, Quota, Convocatória, Atestado, Reserva de Instalação e Lesão. Identifica os estados e as transições, e cruza com os RFs respetivos (RF-08, RF-11, RF-12, RF-17, RF-20).}{}

\subsubsection{Diagramas de Sequência e Atividade}

Os diagramas de sequência foram pedidos para os fluxos críticos: autenticação JWT, registo de utilizador, pagamento de quota, publicação de convocatória e reserva de instalação. O Gemini CLI gerou-os a partir das tabelas de caso de uso correspondentes.

Os diagramas de atividade foram gerados a partir dos mesmos casos de uso, mas com foco no fluxo de controlo interno em vez das interações entre componentes.

\caixaprompt{Gemini CLI}{Olha para a tabela do caso de uso de pagamento no Cap. 4 e gera o diagrama de sequência em PlantUML. Os participantes são o utilizador, o frontend, a API, o módulo financeiro, o PostgreSQL, o gateway de pagamentos e as notificações. Inclui o webhook de confirmação e o que acontece se o gateway falhar.}{}

\caixaprompt{Gemini CLI}{Gera o diagrama de atividade do mesmo caso de uso, focado no fluxo interno (validações, decisões, ações paralelas).}{}

\subsubsection{Casos de Uso e Especificação de API}

As 27 tabelas de caso de uso e a especificação dos sete \textit{endpoints} REST nucleares foram redigidas com Claude Code, partindo dos diagramas e do esquema de requisitos. O modelo identificou que o \textit{endpoint} \texttt{GET /api/athletes/\{id\}/status} era crítico para o \textit{dashboard} do treinador e deveria existir como endpoint dedicado (em vez de exigir que o cliente compusesse o resultado a partir de chamadas a vários módulos) --- uma sugestão que a equipa adotou imediatamente.

\caixaprompt{Claude Code (CLI)}{Olha para as tabelas de caso de uso que já existem no chapters/04\_especificacao\_modelacao.tex e gera as restantes para todos os 23 RF. Usa exatamente o mesmo formato. Mete um parágrafo antes de cada tabela a explicar o caso de uso.}{}

\caixaprompt{Claude Code (CLI)}{Revê os endpoints REST que já estão documentados. Os métodos HTTP, os códigos de resposta e o controlo de acesso estão corretos? Falta algum endpoint crítico?}{}

\subsubsection{Design de Interfaces e Mockups (V0)}

Para o desenho das interfaces das principais páginas da aplicação utilizámos o Vercel v0. A vantagem desta ferramenta é que devolve diretamente código React/JSX com Tailwind CSS, pronto a ser integrado no \textit{frontend}. A equipa descreveu cada um dos quatro \textit{dashboards} principais (Atleta, Treinador, Encarregado de Educação, Secretaria) e o v0 gerou versões iniciais que serviram de base para a implementação.

Foram também gerados \textit{wireframes} estáticos em PlantUML (subsistema Salt) para incluir no relatório como protótipos de baixa fidelidade. Estes wireframes serviram para discussão interna antes de avançar para o código real.

\caixaprompt{v0 (Web)}{Desenha o dashboard do treinador para um clube desportivo. Quero uma sidebar com Plantel, Presenças, Convocatórias, Calendário, Atestados e Relatórios. No topo uns cards com os números importantes (atletas, assiduidade, próximo treino). Depois uma lista de próximos eventos e um painel com alertas tipo atestados caducados ou quotas em atraso. Estilo limpo, azuis e brancos, usa shadcn/ui.}{}

\caixaprompt{v0 (Web)}{Agora o dashboard do encarregado de educação. Precisa de um seletor de filho no topo (pode ter mais que um no clube), um semáforo com o estado do atleta (atestado, lesão, quotas), as próximas convocatórias com botão de confirmar, e a parte financeira com referência MB/MBWay para pagar e histórico de faturas.}{}

\caixaprompt{Gemini (Web)}{Que princípios de UX recomendas para uma aplicação usada por perfis tão diferentes como um atleta de 12 anos, um treinador em campo e um funcionário de secretaria? Há ecrãs que merecem atenção especial?}{}

\subsection{Fase 4 --- Implementação da Aplicação}\label{sec:ia_implementacao}

A fase de implementação foi onde o Claude Code mais brilhou. A capacidade de operar diretamente sobre múltiplos ficheiros C\#, configurações de \texttt{appsettings.json}, ficheiros \texttt{Dockerfile} e \texttt{docker-compose.yml} num único passo permitiu reduzir significativamente o tempo gasto em \textit{boilerplate}. A equipa manteve o controlo das decisões arquiteturais e da revisão dos \textit{diffs}, mas a redação propriamente dita do código foi maioritariamente assistida.

\subsubsection{Estratégia Geral de Geração de Código}

A estratégia adotada foi a seguinte:

\begin{enumerate}
    \item \textbf{Esqueleto manual.} A equipa criou manualmente a estrutura de pastas do projeto, os ficheiros \texttt{.csproj} e o \texttt{docker-compose.yml} inicial, para garantir que a base de configuração estava sob controlo humano.
    \item \textbf{Entidades de domínio com Claude Code.} A partir do diagrama de classes, pedimos ao Claude para gerar as entidades C\# correspondentes em cada módulo, respeitando a divisão por esquema na base de dados.
    \item \textbf{Repositórios e DbContexts com Claude Code.} O padrão Repository genérico (\texttt{IRepository<T>}) foi gerado a partir das entidades, com o ciclo de vida \texttt{Scoped} registado no \texttt{Program.cs}.
    \item \textbf{Controllers com Claude Code.} A partir das tabelas de caso de uso e da especificação de API, pedimos a geração dos \texttt{Controllers} com os atributos \texttt{[Authorize(Roles=...)]} corretos.
    \item \textbf{DataSeeder com Claude Code.} O \texttt{DataSeeder} programático foi inteiramente gerado pelo Claude, a pedido nosso para criar dados realistas e variados.
    \item \textbf{Frontend com Cursor + Claude Code.} A camada de apresentação Next.js foi escrita com auxílio do Cursor (no IDE) e completada com Claude Code em \textit{batch} para tarefas como a função utilitária \texttt{fetchApi} ou a estrutura de \textit{routing} por papel.
\end{enumerate}

\subsubsection{Geração das Entidades de Domínio}

\caixaprompt{Claude Code (CLI)}{Lê o diagrama\_classes.puml e gera as entidades C\# para os 7 módulos em src/Modules/. Cada entidade herda de Entity com Id Guid. Configura TPH para a hierarquia de Utilizador. Cada módulo tem o seu DbContext e o seu esquema PostgreSQL separado (users, sports, clinical, etc.).}{}

\subsubsection{Geração dos Controllers e Lógica de Negócio}

\caixaprompt{Claude Code (CLI)}{Gera os Controllers REST para cada módulo a partir dos casos de uso no Cap. 4. Mete os [Authorize(Roles=...)] corretos conforme a tabela RBAC. No AuthController mete o bloqueio depois de 5 tentativas falhadas.}{}

\caixaprompt{Claude Code (CLI)}{Implementa o JWT no AuthController. Claims: NameIdentifier, Name, Email e Role. Validade de 7 dias, HMAC-SHA256, chave no appsettings.json.}{}

\subsubsection{Implementação do DataSeeder}

\caixaprompt{Claude Code (CLI)}{Cria um DataSeeder em src/CAP.API/Data/ que popule a BD se estiver vazia. Mete funcionários, treinadores, EEs e bastantes atletas distribuídos por equipas. Gera treinos passados com presenças variadas (usa Random), atestados em vários estados (válido, a expirar, expirado) e quotas também (paga, pendente, em atraso). Usa datas relativas a DateTime.UtcNow para ficar sempre realista.}{}

\subsubsection{Camada de Apresentação (Frontend)}

\caixaprompt{Claude Code (CLI)}{Monta a estrutura de routing do frontend Next.js em src/CAP.Frontend/app/. Cada papel (atleta, treinador, encarregado, secretaria) tem a sua pasta com as páginas. Cria também um fetchApi em lib/api.ts que meta automaticamente o JWT no header e redirecione para login se der 401.}{}

\caixaprompt{v0 (Web)}{Faz o dashboard da Secretaria em React com o wireframe do Cap. 4. Precisa de ter as aprovações pendentes em tabs, botão de exportar SAF-T e botão de registar pagamento manual. Usa shadcn/ui e Tailwind.}{}

\subsection{Fase 5 --- Redação do Relatório}\label{sec:ia_relatorio}

Para além das fases anteriores, o LLM foi também utilizado na redação do próprio relatório. Foram três os tipos de utilização principais:

\subsubsection{Redação em Bloco do Texto Explicatório}

Como mencionado na Secção~\ref{sec:ia_requisitos}, a redação do parágrafo explicativo de cada um dos 36 requisitos foi delegada ao Claude Code, com o \textit{prompt} contendo o nome do requisito, a sua origem e as restrições do CAP. Esta abordagem em \textit{batch} poupou horas de trabalho, mas teve a contrapartida de gerar texto demasiado uniforme nos primeiros seis requisitos (RF-01 a RF-06), que tiveram de ser reescritos manualmente pela equipa para introduzir variação de estilo (ver Secção~\ref{sec:ia_validacao}).

\subsubsection{Correção Gramatical e Estilística}

\caixaprompt{Claude Code (CLI)}{Corre os ficheiros todos em chapters/ e corrige gralhas e erros gramaticais. Não mudes o conteúdo, só a qualidade da escrita. Diz-me o que mudaste.}{}

\subsubsection{Validação e Limpeza de Sintaxe LaTeX}

\caixaprompt{Claude Code (CLI)}{Compila o relatório e diz-me todos os warnings que encontras: referências partidas, labels duplicados, overfull boxes, citações por resolver. Corrige os que forem óbvios.}{}

% ============================================================
% SECÇÃO 6.3 — ENGENHARIA DE PROMPTS
% ============================================================

\section{Engenharia de \textit{Prompts}: Boas Práticas e Lições}\label{sec:ia_prompts}

Ao longo do projeto fomos refinando a forma como interagimos com os modelos. Algumas das lições mais úteis ficaram registadas aqui.

\subsection{Boas Práticas Identificadas}

\begin{itemize}
    \item \textbf{Fornecer contexto antes de pedir.} Sempre que abrimos um chat novo, começámos por carregar o enunciado do projeto, a lista de requisitos consolidada ou o ficheiro \texttt{.tex} relevante. Os \textit{prompts} sem contexto produziam respostas genéricas e tipicamente alucinadas.

    \item \textbf{Especificar restrições concretas em vez de termos abstratos.} ``Considera RGPD'' é vago; ``garante que dados clínicos só são acessíveis ao Atleta, ao EE e ao próprio médico, e nunca ao Treinador'' é acionável. A diferença na qualidade da resposta é substancial.

    \item \textbf{Pedir formato explícito.} Tabelas em \LaTeX{} ($|l|X|$), código em Markdown (\texttt{```csharp}), JSON com \textit{schema} indicado. Quando o formato fica ambíguo, o modelo escolhe um genérico que raramente é o que queremos.

    \item \textbf{Iterar em vez de pedir tudo de uma vez.} Tentar fazer ``gera-me todos os RF detalhados de uma só vez'' produz texto uniforme. Pedir um a um, ou em lotes de 3--5 com instrução para variar o estilo, dá melhor resultado.

    \item \textbf{Abrir um chat novo para tarefas independentes.} A memória do chat pode contaminar tarefas distintas com pressupostos da conversa anterior. Para extrair requisitos de forma neutra, por exemplo, abrimos um chat sem histórico.
\end{itemize}

\subsection{Anti-Padrões que Encontrámos}

\begin{itemize}
    \item \textbf{``Faz como achares melhor''.} O modelo escolhe a opção mais comum, que raramente é a mais adequada ao contexto específico. Resulta em respostas plausíveis mas pouco úteis.

    \item \textbf{Delegar a redação de várias secções consecutivas sem variar o \textit{prompt}.} Produz uniformidade estilística e expressões repetidas (``pilar estrutural'', ``fulcral'', ``transversal a toda a dinâmica'').

    \item \textbf{Aceitar a primeira versão sem cruzar com fontes.} O modelo é capaz de inventar APIs, métodos e padrões que não existem. Em código, isto traduz-se em erros de compilação imediatos; em texto técnico, em afirmações falsas que passam despercebidas.

    \item \textbf{Não guardar os \textit{prompts}.} Sem registo, é impossível reproduzir resultados nem aprender com tentativas anteriores. Por isso, todos os \textit{prompts} relevantes deste projeto foram guardados num diário de bordo da equipa --- e estão reproduzidos neste capítulo.
\end{itemize}

% ============================================================
% SECÇÃO 6.4 — VALIDAÇÃO HUMANA E REFLEXÃO CRÍTICA
% ============================================================

\section{Validação Humana e Reflexão Crítica}\label{sec:ia_validacao}

Usar LLMs traz riscos concretos: inventar APIs que não existem, sugerir soluções genéricas que não servem ao nosso caso, gerar texto todo igual e dar uma falsa sensação de qualidade (o texto parece bom mas o conteúdo está errado). Esta secção documenta o que correu mal e como lidámos com isso.

\subsection{O que a IA Falhou ou Inventou}

Ao longo do projeto, encontrámos vários casos em que os modelos produziram conteúdo incorreto ou inadequado ao contexto do CAP. Os mais significativos foram os seguintes:

\begin{itemize}
    \item \textbf{Recomendação inicial de microsserviços.} Na fase de decisão arquitetural, o primeiro \textit{prompt} sem restrições explícitas levou o Gemini a recomendar uma arquitetura de microsserviços --- desproporcionada para uma equipa de 4 pessoas e um prazo de 16 semanas. Foi necessário reforçar o contexto (equipa reduzida, prazo curto, projeto académico) para obter uma sugestão razoável de Monólito Modular.

    \item \textbf{Atributos no modelo de domínio.} Apesar de instruções explícitas, o Gemini insistiu por três iterações em incluir atributos nas entidades do modelo de domínio (que devia conter apenas conceitos e relações). Foi corrigido manualmente.

    \item \textbf{Método de pagamento inexistente.} Numa versão inicial do RF-17 (Pagamento de Quotas), o Gemini incluiu ``pagamento por criptomoedas'' como opção --- algo completamente fora do contexto de um clube desportivo amador português, mas que ele evidentemente tinha visto em \textit{datasets} sobre \textit{fintech}. Removemos.

    \item \textbf{Stakeholder em falta: Encarregado de Educação.} Na primeira identificação de \textit{stakeholders}, o modelo não listou os EE, talvez por não associar imediatamente ``clube desportivo'' a ``atletas menores''. Acrescentámos manualmente.

    \item \textbf{Padrão \textit{fat gateway}.} A primeira versão do diagrama de componentes colocava lógica de negócio no API Gateway, o que é um antipadrão arquitetural. Foi detetado em revisão de equipa e corrigido.

    \item \textbf{API REST com \textit{endpoint} de \textit{login} em falta.} Numa versão intermédia da especificação de API, faltava o \textit{endpoint} \texttt{POST /api/auth/login} --- precisamente o mais crítico do sistema. Foi adicionado depois de o detetarmos no cruzamento com o caso de uso ``Iniciar Sessão''.

    \item \textbf{Tom robótico nas tabelas de RF.} O \textit{batch} de redação dos parágrafos explicativos dos RF gerou texto uniforme demais, com expressões repetidas como ``pilar estrutural'', ``culminar do ciclo'' e ``fulcral''. Foram reescritos pela equipa.
\end{itemize}

\subsection{Processo de Revisão da Equipa}

Para mitigar estes riscos, a equipa adotou um processo de revisão sistemático para todo o conteúdo gerado por IA:

\begin{enumerate}
    \item \textbf{Cruzamento com fontes oficiais.} Todo o código gerado foi cruzado com a documentação das tecnologias (Entity Framework Core, Next.js, PostgreSQL, ASP.NET Core). Toda a sintaxe LaTeX gerada foi compilada e os \textit{warnings} corrigidos.

    \item \textbf{Cruzamento com os requisitos.} Sempre que um artefacto era gerado (UC, diagrama, controller), verificávamos se cumpria os RF/RNF aplicáveis. Quando havia conflito, prevalecia o requisito.

    \item \textbf{Revisão em pares.} Cada elemento da equipa revia o trabalho gerado com auxílio do modelo por outro elemento, antes de ser integrado no relatório ou no código. Esta dupla revisão apanhou várias inconsistências que teriam passado.

    \item \textbf{\textit{Commits} pequenos.} O código foi sendo integrado em \textit{commits} pequenos, para que qualquer alteração indesejada gerada pelo modelo pudesse ser facilmente revertida via \texttt{git}.
\end{enumerate}

\subsection{Estimativa de Impacto}

É difícil quantificar com precisão o impacto da assistência de IA, mas a equipa estima conservadoramente que a sua utilização poupou entre 80 e 120 horas de trabalho equivalente. As tarefas onde o ganho foi maior:

\begin{itemize}
    \item Redação detalhada dos 36 requisitos: \textbf{$\approx$ 30 h poupadas} face à redação manual de cada tabela.
    \item Geração das 27 tabelas de caso de uso: \textbf{$\approx$ 20 h poupadas}.
    \item Geração de \textit{boilerplate} de C\# (DbContexts, Repositories, Controllers): \textbf{$\approx$ 15 h poupadas}.
    \item DataSeeder programático com dados variados: \textbf{$\approx$ 8 h poupadas}.
    \item Correção gramatical e estilística do relatório: \textbf{$\approx$ 10 h poupadas}.
\end{itemize}

Estas horas foram redirecionadas para tarefas onde o juízo humano é insubstituível: validação de decisões arquiteturais com a gerência, revisão crítica do que vinha dos modelos, integração de todas as peças no relatório final e teste manual da aplicação.

\subsection{Conclusão da Reflexão}

No fundo, a nossa experiência mostra que os LLMs funcionam melhor como aceleradores do que como substitutos. As decisões importantes --- arquitetura, divisão em módulos, permissões --- foram todas tomadas pela equipa. O que o LLM fez foi poupar-nos o trabalho mecânico de transformar essas decisões em código, diagramas e texto.

A transparência deste capítulo é tão importante quanto a própria utilização. Ao registar o que pedimos e o que veio dos modelos, qualquer pessoa que leia o relatório consegue distinguir o que foi rascunhado automaticamente do que foi decidido por nós.

%==========================================================================
% FIM DO CAPÍTULO 6
%==========================================================================
